\chapter{实验与系统验证}
\label{ch:experiments}

\section{实验原则}
每次实验都应能回答四个问题：运行了哪一版代码、使用了哪些参数、记录了哪些输入/输出、如何从数据得到结论。实验目录至少保存 metadata、rosbag 或视频、指标、日志摘要和异常说明。

\section{验证层级}
\begin{description}
  \item[单元测试] 验证坐标转换、几何三角测量、A*、配置校验、退化数学和指标计算。
  \item[节点契约测试] 验证 Topic 名称、消息类型、启动参数、真值隔离和包依赖。
  \item[SIL 集成测试] 验证 Stonefish--bridge--估计--控制--推进器闭环。
  \item[HIL 测试] 验证 micro-ROS 串口、MCU 控制核、状态回传和安全退出。
  \item[水池/现场测试] 验证真实传感器、真实外参、任务动作和长期稳定性。
\end{description}

\section{轨迹和控制指标}
评测节点读取估计轨迹和 Ground Truth，计算 ATE/RPE、最大位置误差、运行时长、传感器可用率、退化事件数量和控制环性能。典型定义为
\begin{equation}
  \mathrm{ATE}_{\mathrm{RMSE}}
  = \sqrt{\frac{1}{N}\sum_{i=1}^{N}
  \left\|\mathbf{p}^{est}_i-\mathbf{p}^{gt}_i\right\|^2},
\end{equation}
其中 Ground Truth 只在评测内部对齐原点，不能参与控制。RPE 用于观察局部运动误差；控制评测还应记录频率、周期均值、标准差、最大 jitter 和 deadline miss。

\Needspace{12\baselineskip}
\section{推荐记录 Topic}
\begin{lstlisting}[caption={基线实验的最小 rosbag 主题集合}]
/auv/sensors/imu/data
/auv/sensors/dvl/velocity
/auv/sensors/pressure
/auv/state/odom
/auv/state/twist
/auv/state/health
/auv/control/trajectory
/auv/sim/ground_truth/odom
/auv/evaluation/metrics
/auv/sim/degradation/events
/auv/sim/control_performance
\end{lstlisting}

\section{可复现实验入口}
\path{uv_sim_bringup/launch/experiment.launch.py} 支持选择 world、seed、退化 profile、估计器、运行时长、结果目录和 rosbag 记录。建议实验目录结构如下：
\begin{lstlisting}[caption={实验输出目录约定}]
results/<run-id>/
  metadata.json
  metrics.json
  trajectory.csv
  rosbag/
  logs/
  figures/
\end{lstlisting}

\section{故障注入矩阵}
至少应覆盖 nominal、DVL 丢失、视觉丢失、DVL+视觉同时退化、USBL 离群、控制线程 deadline miss、相机启动失败、MCU heartbeat 超时和任务动作超时。每个场景应规定预期：车辆是否停止、估计器是否降级、任务是否重试、指标是否被标记为无效。

\section{实验报告模板}
实验报告应记录目标、设备和场景、参数 profile、传感器标定版本、软件 commit、seed、数据路径、指标、图表、异常、结论和限制。若结果来自仿真，不得省略与真机差异和 Ground Truth 使用范围。
